\chapter{研究结论与实践启示}
\section{研究框架与待验证结论}
本文围绕中国股票已实现波动率预测，提出由高频测度、动态条件风险关系与边特征消息传递构成的研究框架。与仅比较复杂模型和简单HAR不同，本文将风险汇总、具体关系、动态相关和边际尺度放入统一的贡献归因协议，以检验逐边条件化是否具有额外价值。

在方法层面，本稿已经明确三个结论边界。首先，所选对称边际与t-Copula设定中的上下尾强度具有镜像对应，不能预先作为两个独立信号。其次，原始条件风险量的方向差可能包含接收方波动尺度的作用，标准化与分布形状控制是解释方向的重要步骤。最后，风险汇总在一般图结构中不能决定逐边加权邻居状态，但这种表示差异是否有实际预测收益，必须由有限样本实验确认。

上述内容属于模型定义与数学推导，不替代市场实证。最终结论应在正式结果生成后分别说明：完整模型相对匹配风险汇总与动态相关模型的增量、有效边特征、适用期限和状态，以及计算与数据条件。

\Pending{根据第五章核心表格回填主要实证结论；每条结论注明适用样本、效应大小和不确定性。}

\section{实践启示}
\subsection{风险预测与数据治理}
对实际使用而言，风险预测结果应与数据时段、收益调整口径和失败标志一起解释。只给出模型输出而不说明其输入是否完整，容易把数据缺失造成的平稳误判为真实低风险。公司数据适配中的时间编码、参考价、公司行动版本与历史成分证据，因而不是模型之外可忽略的技术细节。

若逐边模型经样本外比较显示稳定增量，其适当用途是为既有风险预测增加条件关系信息，而不是把高权重边直接命名为市场中的因果传染路径。若风险汇总已经取得相近效果，则较简约的模型在运行成本和维护方面可能更合适；这一取舍应结合实际差值与成本报告判断。

\subsection{资产配置及监管解释的边界}
本文的主目标是方差预测，而非交易收益最大化。将方差预测用于组合权重、VaR或ES，还需额外的收益分布、协方差结构、持有区间及执行约束。没有对应回测证据时，不将预测误差降低写成已经提高净收益、降低回撤或识别系统重要性主体。

对于市场监测，条件风险网络可以作为描述和研究工具，但其解释应保留模型假设、样本代表性和边际尺度限制。固定股票样本中的输入、输出排名不能直接外推为整个市场的风险贡献排序。

\section{研究局限与未来方向}
本研究的首要数据局限是公司历史分钟与辅助数据的版本可得性。若只有当前修订的历史快照，严格的算法时间顺序仍不能完全还原历史线上环境；这一限制需在最终论文中保留。其次，二元关系模型的拼接不保证全维联合分布有效，因而本研究不直接解决完整协方差与组合风险分配问题。

模型方面，标准t-Copula限制了上下尾非对称的表达，浅层图模型则在表示能力与有限样本稳定性之间作出取舍。未来可以在保持同一强基线和时间协议的前提下，比较非对称依赖、更多证券以及不同交易时段的外部效度，而不是仅因模型更复杂就视为自然改进。

在最终研究结论中，应把方法结构所保证的性质、真实样本所支持的结果和仍需验证的应用分别陈述。这样的区分使论文从研究问题到证据形成闭合，也为后续复现和扩展保留清楚的起点。
